\tema{Antes de empezar: símbolos y operaciones básicas}
\paso{1. Una función es una regla de cálculo.}
Si escribimos $F(x)$, $x$ es la entrada y $F(x)$ la salida. Por ejemplo, si $F(x)=2+3x$, entonces $F(0.1)=2+3(0.1)=2+0.3=2.3$. Los paréntesis después del nombre indican evaluación; no significan que se multiplique $F$ por $x$.

\paso{2. Potencias y signos.}
$x^2=x\cdot x$, $x^3=x\cdot x\cdot x$ y $x^0=1$, también en la evaluación del término constante de un polinomio en cero. Para $x=0.1$: $x^2=0.01$ y $x^3=0.001$. Multiplicar dos negativos da positivo: $(-2)(-4.6095)=9.219$. Multiplicar un negativo por un positivo da negativo: $(-6)(0.741)=-4.446$.

\paso{3. Distribuir y agrupar.}
$a(b+c)=ab+ac$. Por ejemplo, $x(1-x)=x-x^2$. Solo se suman directamente términos de igual potencia: $2x^2+3x^2=5x^2$, pero $2x+3x^2$ conserva dos términos.

\paso{4. Logaritmo natural y exponencial.}
$\ln$ significa logaritmo natural, de base $e\simeq2.71828$. Las operaciones $\ln$ y $\exp$ se deshacen mutuamente: si $F=\ln\gamma$, entonces $\gamma=\exp(F)=e^F$. Ejemplo: si $F=-2$, $\gamma=e^{-2}\simeq0.135335$. Un logaritmo negativo corresponde a un coeficiente positivo menor que uno; el coeficiente no es negativo.

\paso{5. Factorial y sumatoria.}
$0!=1$, $1!=1$, $2!=2\cdot1=2$, $3!=3\cdot2\cdot1=6$, $4!=24$. La expresión $\sum_{k=0}^{3}c_kx^k$ significa escribir y sumar exactamente cuatro términos:
\[
c_0x^0+c_1x^1+c_2x^2+c_3x^3=c_0+c_1x+c_2x^2+c_3x^3.
\]
$k$ es un contador. $c_k$ es el coeficiente asociado a la potencia $x^k$.

\paso{6. Derivada, integral y aproximación.}
$F'(x)$ es la derivada: mide cuánto cambia $F$ por un cambio pequeño de $x$. $F''$ es la derivada de $F'$. Una integral definida suma contribuciones acumuladas entre dos límites. Las reglas necesarias se explican antes de aplicarlas. El signo $=$ indica igualdad; $\simeq$ indica un valor aproximado o redondeado.

\paso{7. Nombres que se usarán.}
$f=g^E/(RT)$ es la energía de exceso dividida entre $RT$; $F=\ln\gamma_{\Al}$ es el logaritmo del coeficiente de Al; $H=\ln\gamma_{\Fe}$ es el del Fe. Son funciones diferentes. $P_N$ aproxima $F$ reteniendo hasta grado $N$; $Q_m$ aproxima directamente $\gamma$ hasta grado $m$.

\tema{a) El sistema físico: qué sucede en el metal}
\paso{1. Materiales presentes.}
El hierro líquido es el medio principal. Contiene aluminio y oxígeno disueltos. Los corchetes $[\Al]$ y $[\mathrm O]$ señalan especies disueltas, no necesariamente sustancias puras separadas. Al agregar aluminio puede formarse alúmina sólida:
\[
2[\Al]+3[\mathrm O]\longrightarrow\mathrm{Al}_2\mathrm O_3(s).
\]
El 2 cuenta átomos de Al y el 3 cuenta átomos de O por unidad de alúmina. $(s)$ significa sólido. Disminuye el oxígeno disuelto, pero pueden quedar partículas de alúmina que requieren separación.

\paso{2. Lo que estudia el artículo.}
Estudia equilibrio termodinámico y no idealidad: cómo se relacionan composición, actividad y energía de Gibbs. No calcula cuánto tarda la reacción, cómo se agita el metal o a qué velocidad se retiran las partículas.

\paso{3. Fracciones molares.}
Un mol es una cantidad de sustancia. Si hay $n_{\Al}$ moles de Al y $n_{\Fe}$ de Fe en el modelo binario, el total es $n=n_{\Al}+n_{\Fe}$ y
\[
x_{\Al}=\frac{n_{\Al}}{n},\quad x_{\Fe}=\frac{n_{\Fe}}{n},\quad
x_{\Al}+x_{\Fe}=\frac{n_{\Al}+n_{\Fe}}n=1.
\]
Definimos $x=x_{\Al}$; por tanto $x_{\Fe}=1-x$. Ejemplo: 1 mol de Al más 9 moles de Fe da $n=10$, $x=1/10=0.1$ y $x_{\Fe}=9/10=0.9$. Son fracciones de moles, no de masa.

\paso{4. Actividad y coeficiente de actividad.}
La actividad describe la participación termodinámica de una especie. En la escala utilizada, $a_i=x_i\gamma_i$. Si fuera ideal, $\gamma_i=1$ y $a_i=x_i$. Si $x_i=0.1$ y $\gamma_i=0.05$, entonces $a_i=0.1(0.05)=0.005$. Comparar actividades exige conservar el mismo estado de referencia.

\paso{5. Relación con la desoxidación.}
Para alúmina pura, su actividad es 1. El artículo define la constante de disociación $K_d=a_{\Al}^2a_{\mathrm O}^3$. Dividiendo ambos lados entre $a_{\Al}^2$:
\[
\frac{K_d}{a_{\Al}^2}=a_{\mathrm O}^3.
\]
Tomando raíz cúbica: $a_{\mathrm O}=(K_d/a_{\Al}^2)^{1/3}$. A $T$ fija, $K_d$ es fijo: aumentar $a_{\Al}$ reduce $a_{\mathrm O}$. Para obtener $x_{\mathrm O}=a_{\mathrm O}/\gamma_{\mathrm O}$ hace falta también $\gamma_{\mathrm O}$. El cálculo binario de este trabajo no proporciona ese equilibrio completo.

\tema{a) De tres componentes a dos: hipótesis y unidades}
\paso{1. Qué se simplifica.}
El sistema completo Fe-Al-O es ternario. El artículo propone que, cuando la contribución de O es pequeña, Fe y Al pueden describirse aproximadamente como un binario. Esto es una hipótesis física, anterior a cualquier Taylor. El artículo sitúa esa reducción por encima de aproximadamente 0.07\,\% en masa de Al en las condiciones estudiadas [1, sección 3.2].

\paso{2. Convertir 0.07 por ciento en masa a fracción molar.}
``Por ciento'' significa por cada 100. Tomemos 100 g de mezcla binaria:
\[
m_{\Al}=0.07\ \mathrm g,\qquad m_{\Fe}=100-0.07=99.93\ \mathrm g.
\]
Con masas molares $M_{\Al}=26.98$ g/mol y $M_{\Fe}=55.85$ g/mol, la regla $n=m/M$ da
\[
n_{\Al}=\frac{0.07}{26.98}\simeq0.002594514\ \mathrm{mol},\qquad
n_{\Fe}=\frac{99.93}{55.85}\simeq1.789256938\ \mathrm{mol}.
\]
Sumamos: $n\simeq0.002594514+1.789256938=1.791851452$ mol. Dividimos la cantidad de Al entre el total:
\[
x_{\Al}\simeq\frac{0.002594514}{1.791851452}\simeq0.001447952.
\]
Así se obtiene $x\simeq0.00145$, que no es ni $0.07$ ni $0.0007$. Este cálculo supone solo Fe y Al; no es una conversión exacta de toda composición ternaria.

\paso{3. Temperatura y energía.}
Usamos $T=1873$ K. En grados Celsius, $1873-273.15=1599.85\,{}^\circ$C. La constante es $R=8.314462618$ J/(mol K). Su producto es
\[
RT=(8.314462618)(1873)=15572.986483514\ \mathrm{J/mol}.
\]
Los kelvin se cancelan: $[\mathrm{J/(mol\ K)}][\mathrm K]=\mathrm{J/mol}$. Por eso $g^E/(RT)$ no tiene unidades cuando $g^E$ está en J/mol. $1$ kJ $=1000$ J: para convertir J/mol a kJ/mol dividimos entre 1000.

\paso{4. Qué es energía de exceso.}
$g^E$ es la diferencia entre la energía de Gibbs molar de la solución real modelada y la de una solución ideal de igual composición y mismos estados de referencia. ``Molar'' significa por mol. $g^E<0$ no significa, por sí solo, que la mezcla libere esa cantidad de calor.

\tema{Mapa de procedencia: de dónde sale cada número}
\begin{longtable}{p{0.23\linewidth}p{0.31\linewidth}p{0.36\linewidth}}
\toprule Dato & Valor o expresión & Procedencia y uso\\\midrule\endhead
$A$ & $-4.6095$ & Tabla 1 de [1], coeficiente binario Fe-Al de Redlich-Kister.\\
$B$ & $0.741$ & Tabla 1 de [1], segundo coeficiente binario Fe-Al.\\
Modelo $f$ & $x(1-x)[A+B(1-2x)]$ & Ecuación (16) de [1]. Se toma como punto de partida; no se ajusta aquí a mediciones.\\
$T$ & 1873 K & Condición de estudio de [1].\\
$R$ & 8.314462618 J/(mol K) & Constante física adoptada para convertir energía adimensional a J/mol.\\
$c_0,c_1,c_2,c_3$ derivados & $-3.8685$, $4.773$, $2.0595$, $-2.964$ & Se obtienen de $A,B$ mediante la derivación que se muestra a continuación.\\
$c_2$ publicado & 2.06 & Valor redondeado usado en la ecuación (26) de [1]. Se declara el cambio explícitamente.\\
$x_{\max}$ & 0.3872 & Última composición de Al de la tabla 2 de [1].\\
$M$ & 10\,001 & Elección numérica de este solucionario; no son 10\,001 experimentos.\\
$0.01,0.1,0.2,0.3$ & Composiciones ilustrativas & Elegidas aquí para mostrar cuentas; no se presentan como datos experimentales.\\
$b_0,\ldots,b_{12}$ & Coeficientes de $e^{F_3}$ & Se calculan por una recurrencia derivada aquí. No proceden de otra tabla del artículo.\\\bottomrule
\end{longtable}
\paso{La cadena completa.}
Artículo $\rightarrow A,B,f$ $\rightarrow$ derivación de $F$ $\rightarrow$ Taylor $\rightarrow P_N$ $\rightarrow \gamma_N=e^{P_N}$ $\rightarrow$ errores. En otra rama: $P_N$ $\rightarrow$ Gibbs-Duhem $\rightarrow H_N$ $\rightarrow g_N^E$. Para la extensión infinita: $F_3$ $\rightarrow e^{F_3}$ $\rightarrow b_n$ $\rightarrow Q_m$.

\textbf{No se digitalizan puntos de una imagen ni se inventa una tabla experimental.} Se generan pares $(x,\text{valor calculado})$ evaluando ecuaciones. La tabla 2 sirve como referencia de escala, pero sus últimas cifras no se imponen como si fueran datos exactos de este programa.
